\section{The toy}\label{sec:toy}

\subsection{Machines and the bounded mixture}

Strings are binary; $\eps$ is the empty string, $|x|$ the length of $x$, $x\pre y$ means that $x$ is a prefix
of $y$, and $z_{i:j}$ is the segment $z_i\dots z_j$. Two strings are \emph{incomparable} if neither is a prefix
of the other. $\lambda$ is the fair-coin measure on $\B^{\N}$, and $[p]$ is the set of infinite extensions of
$p$, so $\lambda([p])=2^{-|p|}$.

A \emph{monotone machine} $V$ has a one-way read-only input tape, work tapes, and a one-way write-only output
tape. Given an infinite input $\omega$ it runs forever or halts, and
\begin{itemize}[leftmargin=1.4em,itemsep=1pt,topsep=2pt]
\item $\out_V(\omega)$ is its finite or infinite output;
\item $t_V(\omega,k)\in\N\cup\{\infty\}$ is the step at which the $k$-th output bit is written;
\item $r_V(\omega,k)$ is the number of input bits read by that step.
\end{itemize}
What happens up to step $t_V(\omega,k)$ depends only on the first $r_V(\omega,k)$ bits of $\omega$. Each step
reads at most one input bit, so $r_V(\omega,k)\le t_V(\omega,k)$. If the $k$-th bit is never written,
$t_V(\omega,k)=\infty$ and $r_V(\omega,k)$ plays no role.

A \emph{schedule} is a nondecreasing $T\colon\{1,2,\dots\}\to\N\cup\{\infty\}$, and a length bound is
$L\in\N\cup\{\infty\}$.

\begin{definition}[Bounded mixture]\label{def:mixture}
For a monotone machine $V$, a schedule $T$, a length bound $L$ and a string $x\neq\eps$ let
\[
E_V^{T,L}(x)=\bigl\{\omega:\ x\pre\out_V(\omega),\ \ t_V(\omega,k)\le T(k)\text{ for }1\le k\le|x|,\ \ r_V(\omega,|x|)\le L\bigr\},
\]
let $E_V^{T,L}(\eps)=\B^{\N}$, and $\xi_V^{T,L}(x)=\lambda\bigl(E_V^{T,L}(x)\bigr)$. A string $p$ with $|p|\le L$ and
$[p]\subseteq E_V^{T,L}(x)$ is a \emph{$(T,L)$-program for $x$}; it is a $(T,L)$-program for an infinite $z$
if it is one for every prefix of $z$.
\end{definition}

In words: feed the machine coin flips; $\xi(x)$ is the probability that it prints $x$, on time at every bit,
having read at most $L$ bits. With $T\equiv\infty$ and $L=\infty$ this is the a priori semimeasure $M_V$ of
Solomonoff and Levin~\cite{lv}. We fix one machine $U$ and write $E$, $\xi$ for $E_U^{T,L}$, $\xi_U^{T,L}$ when
the bounds are understood.

\begin{lemma}\label{lem:basic}
\begin{enumerate}[label=(\roman*),leftmargin=2em,itemsep=1pt,topsep=2pt]
\item $E(xb)\subseteq E(x)$ for $b\in\B$, and $E(x)\cap E(x')=\emptyset$ for incomparable $x,x'$. Hence
$\xi(\eps)=1$ and $\xi(x0)+\xi(x1)\le\xi(x)$.
\item If $p$ is a $(T,L)$-program for $x$ then $\xi(x)\ge2^{-|p|}$.
\item If $L<\infty$ and $\xi(x)>0$ then $\xi(x)\ge2^{-L}$.
\end{enumerate}
\end{lemma}
\begin{proof}
(i) The conditions for $xb$ contain those for $x$, because $T$ is imposed bit by bit and $r(\omega,\cdot)$ is
nondecreasing. Incomparable strings cannot both be prefixes of $\out(\omega)$. (ii) is immediate. (iii) Whether
$\omega\in E(x)$ is decided by the first $r(\omega,|x|)\le L$ bits of $\omega$, so $E(x)$ is a union of sets
$[p]$ with $|p|\le L$.
\end{proof}

\begin{definition}[Universality with overhead]\label{def:universal}
$U$ is \emph{universal} if for every monotone machine $V$ there are a word $v$ and a nondecreasing
$\gamma\colon\N\to\N$ such that for all $\omega$ and $k$
\[
\out_U(v\omega)=\out_V(\omega),\qquad r_U(v\omega,k)=|v|+r_V(\omega,k),\qquad t_U(v\omega,k)\le\gamma\bigl(t_V(\omega,k)\bigr).
\]
We call $v$ a \emph{prefix word} and $\gamma$ an \emph{overhead} of $V$, and put $\gamma(\infty)=\infty$;
the same convention applies to the other time bounds introduced below.
\end{definition}

Standard universal machines have this property, with $\gamma(t)=O(t\log t)$ or $O(t^2)$ depending on the tape
model. The results that use universality are
\cref{thm:twin,thm:quine,thm:search,thm:critical,prop:honest,thm:capture,thm:machine} and, in
\cref{sec:stakes}, \cref{prop:stakes-basic}(c), \cref{thm:weighted,prop:aim-cost,prop:twin-stakes} and
\cref{ex:stakes}; the others hold for any machine.

\begin{lemma}[Transfer]\label{lem:transfer}
Let $V$ have prefix word $v$ and overhead $\gamma$. If $\gamma(T'(k))\le T(k)$ for all $k\le|x|$, then
$\xi_U^{T,L}(x)\ge2^{-|v|}\,\xi_V^{T',L-|v|}(x)$.
\end{lemma}
\begin{proof}
If $\omega\in E_V^{T',L-|v|}(x)$ then $v\omega\in E_U^{T,L}(x)$, and $\omega\mapsto v\omega$ multiplies
$\lambda$ by $2^{-|v|}$.
\end{proof}

\subsection{The predictor}

Fix a block length $m\ge1$. For a \emph{context} $x$ let $Z_m(x)=\sum_{y\in\B^m}\xi(xy)$. When $Z_m(x)>0$ the
predictor's distribution over continuations is
\[
\pi_m(y\mid x)=\frac{\xi(xy)}{Z_m(x)}\qquad(y\in\B^m).
\]
The \emph{sampling predictor} returns $Y\sim\pi_m(\cdot\mid x)$; the \emph{MAP predictor} returns a most
probable $y$. Since $Z_m(x)\le\xi(x)$ we have $\pi_m(y\mid x)\ge\xi(xy)/\xi(x)$. With finite $T$ and $L$ the
predictor is computable: it runs the at most $2^{L+1}$ programs for $T(|x|+m)$ steps each.

\subsection{Worlds, handovers, and the hypothesis}

\begin{definition}[World, handover]\label{def:world}
A \emph{world} is an input $\omega$ for $U$; its \emph{transcript} is $\out_U(\omega)$. A \emph{test} is a
decidable set $G$ of nonempty strings; a \emph{first hit} of $G$ is an $x\in G$ no proper prefix of which is in $G$.
The first hits form a prefix-free set. A world \emph{hands over at $x$} if $x$ is the first hit of $G$ among
the prefixes of its transcript and $\omega\in E(x)$. It then receives a block $Y\in\B^m$ from the predictor,
computed from the context $x$, and continues by \emph{running} $Y$, that is, as $U$ on the input $Y0^\infty$.
If $\omega\in E(xd)$ for some $d\in\B^m$, this $d$ is the world's \emph{own continuation}: the next $m$ bits of
its transcript had it not handed over.
\end{definition}

The test with $G=\{x:|x|=n\}$ is ``hand over after $n$ bits''. A test may also look for a marker in the
transcript. Both of H\"anni's suggestions fit: in the second the world is a program that calls the predictor;
in the first, the transcript is the sensor's record and the handover is the deployment (\cref{sec:capture}).
Under this definition every world whose transcript passes the test within the bounds hands over, so whether
a handover takes place can be read off the data. The argument, however, is about predictions that are
actually \emph{used}, and in H\"anni's second formulation only \emph{some} programs call the predictor. The
next definition lets the worlds that use the prediction be any recognisable set of programs.

\begin{definition}[Critical use]\label{def:call}
A \emph{call rule} is a set $\mathcal C$ of pairs $(p,x_p)$ of strings that can be enumerated, one pair for
each \emph{calling program} $p$, such that the calling programs form a prefix-free set and $U$, on every input
beginning with $p$, prints $x_p$ having read no more than $p$. A world that begins with a calling program $p$ \emph{makes a critical use of the
predictor at the context $x_p$}: it receives a block from the predictor, computed from $x_p$, and runs it. Let
\[
C(x)=\sum_{p:\,x_p=x}2^{-|p|},\qquad H_{\mathcal C}=\sum_{p}2^{-|p|},\qquad \kappa(x)=\frac{C(x)}{M_U(x)}
\]
be the prior mass of the worlds that make a critical use at $x$, the prior probability that a world makes a
critical use at all, and the share of the posterior at $x$, without resource bounds, held by worlds that make
one there.
\end{definition}

A test $G$, taken without resource bounds, is the special case in which the calling programs are the strings
$p$ such that $U$ prints a first hit on every input beginning with $p$ and has read exactly $p$ when the hit
is complete: then $C(x)=M_U(x)$ and $\kappa(x)=1$ at first hits, and $H_{\mathcal C}$ is the probability that a
transcript passes the test. H\"anni's second suggestion fits a test only if every program with the given
output calls the predictor. In general a call rule singles out, among the programs that print $x$, those in which the
prediction is actually used; the world is the broader computation around the call. The predictor sees only $x$.

\begin{definition}[Self-reproducing continuation]\label{def:sr}
For a context $x$ let
\[
\SR_m(x)=\bigl\{y\in\B^m:\ y0^\infty\in E(xy)\bigr\}.
\]
\end{definition}

So $y\in\SR_m(x)$ if $U$, run on $y$ padded with zeros, prints $x$ and then $y$ itself, within the bounds.
A world that hands over at a first hit $x$ and receives $y\in\SR_m(x)$ becomes a world whose transcript again
starts with $x$, so under the same test it hands over at $x$ again, and its own continuation is the very $y$
it is running.
A program for $xy$ with $y\in\SR_m(x)$ is a hypothesis that predicts the context and then supplies code that
recreates the situation with itself in the world's place. This is our reading of ``quining oneself out''.

\begin{definition}[Malign prior hypotheses for the toy]\label{def:mph}
\emph{Critical-use form at level $\theta$.} Fix a call rule (or a test) with $H_{\mathcal C}>0$ and, for
each context $x$, a set $A(x)\subseteq\B^m$ of blocks that count as a takeover there. Let a world be drawn from
the prior conditioned on making a critical use, let $X$ be the context of the use and $Y$ the block the
sampling predictor returns, which we assume to be defined at every context.
The hypothesis is that $\Pr[Y\in A(X)]\ge\theta$. With $A(x)=\SR_m(x)$ it says that a share $\theta$ of the
critical uses end with the world running a program that got itself run through the prediction.

\emph{Forms for one world.} Let a world make a critical use at $x$ with own continuation $d$, defined as in
\cref{def:world}.
\begin{itemize}[leftmargin=1.4em,itemsep=1pt,topsep=2pt]
\item \emph{Majority form:} $\pi_m(d\mid x)<\tfrac12$.
\item \emph{Payload form at level $\theta$,} for a set $A\subseteq\B^m$ with $d\notin A$: $\pi_m(A\mid x)\ge\theta$.
\item \emph{Self-reproduction form at level $\theta$:} the payload form with $A=\SR_m(x)\setminus\{d\}$.
\end{itemize}
\end{definition}

The critical-use form is the claim of the argument as stated in \cref{sec:question}: among the predictions
that are actually used in a critical way, a lot end in a takeover through the prediction channel. It is a
statement about the population of uses, and \cref{thm:critical} bounds it from below. It does not exclude
worlds whose own continuation lies in $A(X)$; where the users are themselves such programs it holds for that
reason (\cref{thm:calibration}). The forms for one world say how
much of the posterior at that world's context is off its own continuation and where it sits. Once $U$, $T$,
$L$, $m$ and the world together with the context of its use are fixed, with $T,L<\infty$, these are statements
about finitely many dyadic rationals and can be decided by running the predictor and, for $\SR_m(x)$, by
running $U$ on each candidate $y$ for $T(|x|+m)$ steps. They are necessary for calling the prior malign at
this world, not sufficient: they hold just as well when the world is merely hard to predict
(\cref{thm:calibration}). What one wants to know is how large all these quantities are forced to be, how large
they can be, and through which programs. The rest of the note bounds them.
